\section{符号说明}
\label{sec:4}

\begin{table}[H]
\centering\footnotesize
\caption{主要符号及含义（含单位与角色）}
\label{tab:symbol}
\begin{tabularx}{\textwidth}{@{}p{2.5cm}X p{1.8cm}p{2.5cm}@{}}
\toprule
符号 & 含义 & 单位 & 角色 \\
\midrule
$N$，$D$ & 模型参数量、训练 token 数 & 个，token & Q1/Q2 输入；Q3 决策变量 \\
$C$ & 总训练算力预算 & FLOPs & 给定参数（Q3） \\
$Q(x)$，$Q_i$ & 文档主质量分、训练域 $i$ 质量分 & $[0,1]$ & 估计量（Q1） \\
$Q(\bp)$，$Q_{\rm B}$ & 混合语料质量、B 组质量刻度；主口径二者相同 & 无量纲 & Q1 输出、Q2 输入 \\
$Q_0$ & 推荐配比下的提质起点，$0.6624116679870492$ & 无量纲 & 接口参数（Q2/Q3） \\
$\bp=(p_i)_{i=1}^{17}$ & 17 域训练份额，非负且总和为 1 & 无量纲 & 决策变量（Q1） \\
$z_j(x)$，$w_j$ & 第 $j$ 项标准化指标、冻结的 CRITIC 权重 & $[0,1]$ / 无量纲 & 数据变换、估计参数 \\
$\delta$ & Huber 中心阈值，$0.4034907036619329$ & $[0,1]$ & 由 A1 确定的尺度规则 \\
$B(x)$，$u_s(x)$ & 六来源分歧度、来源 $s$ 的 A1 参照分位 & $[0,1]$ & 诊断量 \\
$L_v$，$L$ & 验证域 $v$ 的损失、等权综合损失 & nats & 观测量或预测目标 \\
$T_{iv}$ & 均匀配比附近域 $i$ 份额对域 $v$ 对数损失的局部效应 & 无量纲 & 估计参数（Q1） \\
$E,A,B,\alpha,\beta$ & 经典标度律地板、幅度与指数 & 依公式 & 估计参数（Q2，B1） \\
$k_0,\kappa_N,\kappa_D$ & 质量地板系数与两项质量幂次 & nats，无量纲 & 估计参数（Q2，B6） \\
$h_N,h_D$ & 主式质量乘子 $Q^{-\kappa_N}$、$Q^{-\kappa_D}$ & 无量纲 & 函数（Q2） \\
$s_{\rm red}$，$\tilde\varphi(\bp)$ & 配比强度（1M 训练定为 $1.4996207320397785$）、配比偏离 & 无量纲 & 参数与由 Q1 给定的量 \\
$m_*$ & 固定 $\bp^*$ 的配比乘子 $\exp\{s_{\rm red}\tilde\varphi(\bp^*)\}\approx0.8989423$ & 无量纲 & Q3 的给定常数 \\
$g(Q)$，$L_{ctx}$，$\eta$ & 单位 token 提质成本、上下文长度、注意力成本系数 $2\times10^{-4}$ & FLOPs/token，token，— & 给定函数与情景参数 \\
$S$，$t$ & Benchmark 综合得分、模型发布日期 & 分，年 & 响应量、时间变量（Q4） \\
$L_{ctx}^{\rm crit}$ & 注意力与训练成本相等的上下文长度，$6/\eta=30000$ & token & 解析结论 \\
\bottomrule
\end{tabularx}
\end{table}

局部符号在首次出现处说明。
